\section*{Part II: Three Statements, One Enterprise}

\subsection*{8. Accounting Is the Language of Obligations}

Accounting is often treated as a compliance function performed after entrepreneurial decisions. In an excellent company it also provides the map of what has been earned, what is owned, what is owed, and how cash changes. These questions overlap but they are not synonymous.

\begin{itemize}
\item \textbf{The income statement} summarizes performance over a period.
\item \textbf{The balance sheet}, or statement of financial position, summarizes assets, liabilities, and equity at a specific moment in time.
\item \textbf{The cash-flow statement} explains movements in cash over a period, commonly grouped into operating, investing, and financing activities.
\end{itemize}
A complete set of statutory financial statements includes further statements, disclosures, and notes; these 3 are the core decision lenses emphasized here.

Standards such as IAS 7 and the Indian Ind AS 7 explain why cash-flow information is distinct from profit information. Revenue recognition under IFRS 15 and its Indian analogue Ind AS 115 focuses on transferring promised goods or services, not on the mere arrival of cash.

\subsection*{9. Income Statement: Did the Period Create Accounting Profit?}

Start with revenue earned during the period. Deduct the expenses recognized for the activities and obligations associated with the period, applying the actual accounting policy. For an education venture these can include instructional delivery, grading, employee compensation, acquisition costs expensed as applicable, platform costs, administration, depreciation, and financing expenses.

A simplified structure is:
\begin{equation}
\begin{aligned}
\text{Net Profit} &= \text{Revenue} - \text{Operating Expenses} \\
&\quad - \text{Financing Costs} - \text{Taxes},
\end{aligned}
\end{equation}
with classification and presentation dependent on applicable standards. Gross margin, contribution, EBITDA, and operating profit are useful intermediate measures, but do not use them interchangeably. Contribution is especially valuable for managerial decisions and is not a standard substitute for a statutory profit measure.

Profit provides evidence about economic performance under recognition and measurement rules. It does not say when receivables will be collected or bills paid. Depreciation reduces reported profit without a cash payment in that period; purchasing an asset consumes cash even when the full purchase amount is not immediately an expense. Both treatments have a reason. Neither permits the founder to ignore the other.

\subsection*{10. Balance Sheet: Where Did the Resources Come From and Where Are They?}

The fundamental identity is:
\begin{equation}
\boxed{\text{Assets} = \text{Liabilities} + \text{Equity}}.
\end{equation}
Assets include cash, customer receivables, inventory, and equipment subject to accounting definitions. Liabilities include supplier payables, customer-service obligations paid in advance, loans, and other amounts owed. Equity represents owners' residual interest after liabilities, including capital invested and accumulated retained profits or losses.

A balance sheet can reveal apparent prosperity and hidden fragility. Consider 2 companies each reporting ₹10 lakh in assets:
\begin{itemize}
\item The 1st has ₹8 lakh of cash and minimal obligations.
\item The 2nd has ₹8 lakh of customer receivables, most overdue, and a large debt payment due this week.
\end{itemize}
The same asset total is not the same liquidity.

Examine \emph{composition}, not just size: how much cash is accessible, which receivables are collectible and when, what inventory may be obsolete, whether equipment can actually be sold, how much unearned service remains, and which liabilities have fixed due dates. The balance sheet is the stock of claims left behind by past decisions.

\subsection*{11. Cash-Flow Statement: Where Did the Cash Come From and Go?}

Under the commonly used IAS 7 framework:
\begin{itemize}
\item \textbf{Operating cash flows} relate to principal revenue-producing activities;
\item \textbf{Investing cash flows} concern acquisitions or disposals of long-lived assets and investments; and
\item \textbf{Financing cash flows} concern changes in contributed equity and borrowings.
\end{itemize}
The applicable classification details should be followed, including treatment of interest and dividends.

The cash-flow statement can be presented using a direct view of actual operating receipts and payments or an indirect reconciliation from an appropriate profit subtotal to operating cash flows. The direct view is often easiest for an early founder planning next week's payments. The indirect view is essential for explaining \emph{why profit is not cash}.

A founder's spreadsheet need not look exactly like the statutory statement to be useful. It must, however, reconcile to actual bank balances and preserve the classification of cash receipts that are loans or owner capital rather than sales. Financing that keeps the lights on is not customer traction.
